\begin{proofbox}
	\textbf{\textcolor{CProof}{方法一（首选速推：椭圆定义＋勾股定理平方差）}}

	\textbf{适用与优势：} 设 $P(x,y)$ 是椭圆上任意一点，$F_1(-c,0)$、$F_2(c,0)$ 是左右焦点，$e=c/a$。本法只用椭圆第一定义、勾股定理、平方差公式与二元一次方程组；\textbf{不必将 $y^2$ 代入椭圆标准方程，也不必配方开根号}。

	\begin{description}[style=nextline, leftmargin=2.8em, labelsep=0.5em, itemsep=0.55em]
		\item[\textbf{第一步（先求和：椭圆定义）：}]
		      记 $r_1=PF_1$，$r_2=PF_2$。因为 $P$ 在椭圆上，所以
		      \[
		          r_1+r_2=2a.
		      \]
		      \par\textbf{理由：} 椭圆上任意一点到两个焦点的距离之和为常数 $2a$。

		\item[\textbf{第二步（求平方差：两次勾股相减）：}]
		      分别对 $PF_1$、$PF_2$ 所在的直角三角形使用勾股定理（等价于两点距离公式）：
		      \[
		          r_1^2=(x+c)^2+y^2,\qquad
		          r_2^2=(x-c)^2+y^2.
		      \]
		      两式相减，$y^2$ 直接消去：
		      \[
		          \begin{aligned}
		              r_1^2-r_2^2
		              &=(x+c)^2-(x-c)^2\\
		              &=\bigl[(x+c)-(x-c)\bigr]\bigl[(x+c)+(x-c)\bigr]\\
		              &=(2c)(2x)=4cx.
		          \end{aligned}
		      \]
		      \par\textbf{理由：} 两个直角三角形的竖直直角边平方相同，消去后只需使用平方差公式。

		\item[\textbf{第三步（平方差除以和，立即得到差）：}]
		      因为 $r_1+r_2=2a>0$，由平方差公式可得
		      \[
		          \begin{aligned}
		              r_1-r_2
		              &=\frac{r_1^2-r_2^2}{r_1+r_2}\\
		              &=\frac{4cx}{2a}\\
		              &=2\frac ca x=2ex.
		          \end{aligned}
		      \]
		      \par\textbf{理由：} 已知两数的和及平方差，用“平方差除以和”得到两数的差；这一过程不涉及开方取正负。

		\item[\textbf{第四步（和差各一半，解出两个焦半径）：}]
		      \[
		          \begin{cases}
		              r_1+r_2=2a,\\
		              r_1-r_2=2ex.
		          \end{cases}
		      \]
		      两式相加、相减并除以 $2$，得
		      \[
		          \boxed{PF_1=a+ex,\qquad PF_2=a-ex.}
		      \]
	\end{description}

	\textbf{考场三行速推：}
	\[
		\begin{aligned}
		    &PF_1+PF_2=2a,\\
		    &PF_1-PF_2=\frac{PF_1^2-PF_2^2}{2a}
		      =\frac{(x+c)^2-(x-c)^2}{2a}=2ex,\\
		    &\therefore\quad PF_1=a+ex,\qquad PF_2=a-ex.
		\end{aligned}
	\]
	\textbf{速记：} 先求和，再作平方差，最后和差各一半。

	\medskip
	\textbf{\textcolor{CProof}{方法二（原有代数证明：代入椭圆方程并配方）}}

	\textbf{\textcolor{CProof}{思路提示}}

	利用两点距离公式表达 $PF_1^2$，借助椭圆方程消去 $y^2$，再利用 $a,b,c,e$ 的关系配方。最后开方时要说明为什么取正号。

	\medskip
	\textbf{\textcolor{CProof}{正式推导}}
	\begin{description}[style=nextline, leftmargin=2.8em, labelsep=0.5em, itemsep=0.55em]
		\item[\textbf{步骤一（写出距离平方）：}]
		      \[
			      PF_1^2 = (x+c)^2 + y^2.
		      \]
		      \par\textbf{理由：} 两点间距离公式
		      \[
			      PF_1=\sqrt{(x-(-c))^2+(y-0)^2},
		      \]
		      平方后得到上式。

		\item[\textbf{步骤二（消去 $y^2$）：}]
		      由椭圆方程得
		      \[
			      y^2=b^2\left(1-\frac{x^2}{a^2}\right).
		      \]
		      代入得
		      \[
			      PF_1^2=(x+c)^2+b^2\left(1-\frac{x^2}{a^2}\right).
		      \]
		      \par\textbf{理由：} 由
		      \[
			      \frac{x^2}{a^2}+\frac{y^2}{b^2}=1
		      \]
		      解出 $y^2$。

		\item[\textbf{步骤三（利用 $b^2=a^2-c^2$ 展开）：}]
		      \[
			      \begin{aligned}
				      PF_1^2
				       & =x^2+2cx+c^2+(a^2-c^2)\left(1-\frac{x^2}{a^2}\right) \\
				       & =x^2+2cx+c^2+a^2-c^2-\frac{a^2-c^2}{a^2}x^2.
			      \end{aligned}
		      \]
		      \par\textbf{理由：} 将 $b^2$ 用 $a^2-c^2$ 代换，再展开合并。

		\item[\textbf{步骤四（合并 $x^2$ 项）：}]
		      \[
			      \begin{aligned}
				      PF_1^2
				       & =\left(1-\frac{a^2-c^2}{a^2}\right)x^2+2cx+a^2 \\
				       & =\frac{c^2}{a^2}x^2+2cx+a^2.
			      \end{aligned}
		      \]
		      \par\textbf{理由：}
		      \[
			      1-\frac{a^2-c^2}{a^2}
			      =\frac{a^2-(a^2-c^2)}{a^2}
			      =\frac{c^2}{a^2}.
		      \]

		\item[\textbf{步骤五（配方并开方）：}]
		      \[
			      PF_1^2=\left(\frac{c}{a}x+a\right)^2.
		      \]
		      因为点 $P$ 在椭圆上，所以
		      \[
			      -a\le x\le a.
		      \]
		      又因为
		      \[
			      0<e=\frac{c}{a}<1,
		      \]
		      所以
		      \[
			      a+ex\ge a-ea=a-c>0.
		      \]
		      因此
		      \[
			      PF_1=\sqrt{(a+ex)^2}=a+ex.
		      \]
		      \par\textbf{理由：} 配方为
		      \[
			      \frac{c^2}{a^2}x^2+2cx+a^2
			      =\left(\frac{c}{a}x+a\right)^2
			      =(a+ex)^2.
		      \]
		      由于 $PF_1$ 是线段长度，且 $a+ex>0$，所以开方后取正号。

		\item[\textbf{步骤六（对应右焦点）：}]
		      同理，
		      \[
			      PF_2^2=(x-c)^2+y^2.
		      \]
		      代入椭圆方程并化简可得
		      \[
			      PF_2^2=(a-ex)^2.
		      \]
		      又因为
		      \[
			      a-ex\ge a-ea=a-c>0,
		      \]
		      所以
		      \[
			      PF_2=a-ex.
		      \]
		      \par\textbf{理由：} 对右焦点 $F_2(c,0)$ 进行同样推导即可。开方时同样要利用 $-a\le x\le a$ 与 $0<e<1$，保证 $a-ex>0$。
	\end{description}

	\medskip
	\textbf{\textcolor{CProof}{结论回扣}}

	因此，对于焦点在 $x$ 轴上的标准椭圆
	\[
		\frac{x^2}{a^2}+\frac{y^2}{b^2}=1\quad (a>b>0),
	\]
	有
	\[
		PF_1=a+ex,\qquad PF_2=a-ex.
	\]
	由此还可得到
	\[
		PF_1+PF_2=2a,\qquad PF_1-PF_2=2ex.
	\]
	其中 $PF_1-PF_2$ 是有向差。
\end{proofbox}